\section{Background and Related Work}
\label{sec:background}

\subsection{Why LLM batching is different}

Conventional request-level batching waits for several inputs, executes one
tensor program, and releases the batch when all outputs are ready. Generative
inference has two asymmetric phases. Prefill processes many prompt tokens in
parallel and is typically compute intensive; decode processes one new token per
active sequence and is often memory-bandwidth limited. Request lengths vary,
so static batches suffer from head-of-line blocking. Orca's iteration-level
scheduler allows the active set to change between decode iterations
\citep{yu2022orca}. PagedAttention reduces fragmentation in the dynamically
growing KV cache and enables \vllm{} to sustain larger active sets
\citep{kwon2023pagedattention}.

Recent systems optimize this inner schedule. Sarathi-Serve uses chunked
prefills and stall-free schedules to reduce prefill--decode interference
\citep{agrawal2024sarathiserve}. DeepSpeed-FastGen's Dynamic SplitFuse also
decomposes prompts and generations to balance throughput and latency
\citep{holmes2024fastgen}. DistServe separates prefill and decode resources to
optimize goodput under distinct latency objectives \citep{zhong2024distserve}.
These systems have detailed visibility into token-level execution. \adip{}
does not: it treats \vllm{} as a black-box HTTP backend and asks what an outer
layer can accomplish with request-level signals alone.

\subsection{Layered prediction serving}

Clipper demonstrated that a serving layer can batch calls to otherwise
unmodified ML frameworks and adapt batch size to a latency objective
\citep{crankshaw2017clipper}. That result motivates gateway batching, but the
backend models in classical prediction serving usually execute a bounded
forward pass. An LLM backend is itself a stateful scheduler. Outer batching
therefore changes both transport granularity and the arrival process observed
by the inner scheduler.

The workload also matters. BurstGPT shows that real LLM services exhibit
time-varying concurrency, response lengths, and failures, and warns that
scheduling conclusions may not generalize across traces
\citep{wang2024burstgpt}. Our primary experiment deliberately uses a
deterministic constant-rate workload. This sacrifices trace realism to isolate
the mechanism; burst sensitivity is left as future validation rather than
being implied by the constant-rate result.

\subsection{Window formation versus backlog formation}
\label{sec:mechanism-model}

Let $\lambda$ be the request arrival rate and $w$ the gateway's intentional
wait window. Under a Poisson approximation, useful for intuition even though
our primary arrivals are deterministic,
\begin{equation}
  \mathbb{E}[N_w] = \lambda w, \qquad
  P(N_w \ge 1) = 1-e^{-\lambda w},
  \label{eq:window}
\end{equation}
where $N_w$ is the number of companion arrivals during the window. At the
largest tested rate, 1.8 requests/s, a 1\,ms window yields
$\mathbb{E}[N_w]=0.0018$ and $P(N_w\ge1)\approx0.18\%$. Even a 20\,ms cap
yields only 0.036 expected companions and a 3.54\% probability. Intentional
micro-waiting alone should therefore produce almost exclusively singleton
batches in this regime.

Now let $S$ be the time for which a serial gateway worker awaits one backend
call. Requests continue to arrive during that interval, creating approximately
$\lambda S$ queued companions. When the worker returns and drains the ingress
queue, it can dispatch a multi-request backend call even if $w$ is nearly zero.
This \emph{backlog coalescing} is qualitatively different from
Eq.~\ref{eq:window}: its control variable is service occupancy, not the
micro-wait. It also exposes the cost of a serial outer scheduler. A request
that arrives behind an in-flight call cannot reach \vllm{} until the outer
worker returns, even though \vllm{} could otherwise admit it into continuous
batching.

Our evaluation tests which mechanism dominates by jointly examining observed
batch size, gateway queue time, backend time, achieved throughput, and tail
latency. This mechanism-level accounting is the paper's main distinction from
a benchmark that merely ranks configurations.
